\documentclass[11pt]{article}
\usepackage{amsmath,amssymb,amsthm,amsfonts}
\usepackage{geometry}
\geometry{margin=1in}
\usepackage{algorithm}
\usepackage{algpseudocode}

\newtheorem{theorem}{Theorem}
\newtheorem{lemma}{Lemma}
\newtheorem{definition}{Definition}
\newtheorem{assumption}{Assumption}
\newtheorem{corollary}{Corollary}

\title{Convergence Analysis of Nesterov's Accelerated Gradient Method\\ for Strongly Convex and Smooth Functions}
\author{}
\date{}

\begin{document}
\maketitle

%==========================================================
\section*{Algorithm Name}
\textbf{Nesterov's Accelerated Gradient Method (NAG) with constant momentum for $\mu$-strongly convex, $L$-smooth functions.}

%==========================================================
\section*{Mathematical Formulation}

The algorithm maintains two interleaved sequences $\{x_k\}$ and $\{y_k\}$. Given the current iterate $x_k$ and previous iterate $x_{k-1}$, the updates are:
\begin{align}
y_k &= x_k + \beta_k\,(x_k - x_{k-1}), \label{eq:extrap}\\
x_{k+1} &= y_k - \alpha\,\nabla f(y_k). \label{eq:grad}
\end{align}

\paragraph{Hyperparameters.}
\begin{itemize}
\item Step size: $\displaystyle \alpha = \frac{1}{L}$.
\item Momentum coefficient (constant in $k$):
\[
\beta_k \equiv \beta = \frac{\sqrt{L}-\sqrt{\mu}}{\sqrt{L}+\sqrt{\mu}}
= \frac{\sqrt{\kappa}-1}{\sqrt{\kappa}+1},
\qquad \kappa := \frac{L}{\mu}\ge 1 .
\]
\end{itemize}
Here $L$ is the smoothness (Lipschitz-gradient) constant and $\mu$ is the strong-convexity constant. Initialization sets $x_{-1}=x_0$, so that $y_0=x_0$.

%==========================================================
\section*{Pseudocode}

\begin{algorithm}[H]
\caption{Nesterov Accelerated Gradient (strongly convex case)}
\begin{algorithmic}[1]
\State \textbf{Input:} $f$, constants $L\ge\mu>0$, initial point $x_0$, tolerance $\varepsilon>0$, max iterations $K$
\State \textbf{Initialize:} $x_{-1}\gets x_0$
\State $\alpha \gets 1/L$
\State $\beta \gets \dfrac{\sqrt{L}-\sqrt{\mu}}{\sqrt{L}+\sqrt{\mu}}$
\For{$k = 0,1,2,\dots,K$}
    \State $y_k \gets x_k + \beta\,(x_k - x_{k-1})$ \Comment{extrapolation}
    \State $g_k \gets \nabla f(y_k)$ \Comment{gradient at extrapolated point}
    \State $x_{k+1} \gets y_k - \alpha\, g_k$ \Comment{gradient step}
    \If{$\|\nabla f(x_{k+1})\| \le \varepsilon$} \textbf{break} \Comment{termination}
    \EndIf
\EndFor
\State \textbf{Output:} $x_{k+1}$
\end{algorithmic}
\end{algorithm}

%==========================================================
\section*{Dry Run Example}

Consider $f(w)=(w-3)^2$, so $\nabla f(w) = 2(w-3)$.
This is $\mu$-strongly convex and $L$-smooth with $\mu = L = 2$ (since $f''=2$).
Then $\kappa = L/\mu = 1$ and
\[
\beta = \frac{\sqrt{2}-\sqrt{2}}{\sqrt{2}+\sqrt{2}} = 0 .
\]
The problem specifies the experimental step size $\eta=\alpha=0.1$. With $\beta=0$ the extrapolation vanishes ($y_k=x_k$), and the method reduces to plain gradient descent $x_{k+1}=x_k-0.1\nabla f(x_k)$.

Initialization: $x_0 = 0$, $x_{-1}=0$.

\paragraph{Iteration $k=0$.}
\begin{align*}
y_0 &= x_0 + 0\cdot(x_0-x_{-1}) = 0,\\
g_0 &= \nabla f(0) = 2(0-3) = -6,\\
x_1 &= y_0 - 0.1\cdot(-6) = 0 + 0.6 = 0.6,\\
f(x_1) &= (0.6-3)^2 = (-2.4)^2 = 5.76 .
\end{align*}

\paragraph{Iteration $k=1$.}
\begin{align*}
y_1 &= x_1 + 0\cdot(x_1-x_0) = 0.6,\\
g_1 &= \nabla f(0.6) = 2(0.6-3) = -4.8,\\
x_2 &= 0.6 - 0.1\cdot(-4.8) = 0.6 + 0.48 = 1.08,\\
f(x_2) &= (1.08-3)^2 = (-1.92)^2 = 3.6864 .
\end{align*}

\paragraph{Iteration $k=2$.}
\begin{align*}
y_2 &= x_2 + 0\cdot(x_2-x_1) = 1.08,\\
g_2 &= \nabla f(1.08) = 2(1.08-3) = -3.84,\\
x_3 &= 1.08 - 0.1\cdot(-3.84) = 1.08 + 0.384 = 1.464,\\
f(x_3) &= (1.464-3)^2 = (-1.536)^2 = 2.359296 .
\end{align*}

The objective values $5.76 \to 3.6864 \to 2.3593$ decrease monotonically toward the minimum $f(3)=0$. (With $\kappa>1$, $\beta>0$ would introduce nonzero momentum and faster acceleration.)

%==========================================================
\section*{Assumptions}

\begin{assumption}[Differentiability]\label{as:diff}
$f:\mathbb{R}^d\to\mathbb{R}$ is continuously differentiable.
\end{assumption}

\begin{assumption}[$L$-smoothness]\label{as:smooth}
$\nabla f$ is $L$-Lipschitz: for all $x,y$,
\[
\|\nabla f(x)-\nabla f(y)\| \le L\,\|x-y\|.
\]
Equivalently, $f(y) \le f(x) + \langle \nabla f(x), y-x\rangle + \frac{L}{2}\|y-x\|^2$.
\end{assumption}

\begin{assumption}[$\mu$-strong convexity]\label{as:sc}
There exists $\mu>0$ such that for all $x,y$,
\[
f(y) \ge f(x) + \langle \nabla f(x), y-x\rangle + \frac{\mu}{2}\|y-x\|^2.
\]
\end{assumption}

We denote by $x^\star$ the unique minimizer, $\nabla f(x^\star)=0$, and $\kappa = L/\mu \ge 1$.

%==========================================================
\section*{Main Convergence Theorem}

\begin{theorem}[Linear (geometric) convergence]\label{thm:main}
Under Assumptions~\ref{as:diff}--\ref{as:sc}, with $\alpha = 1/L$ and
$\beta = \frac{\sqrt{\kappa}-1}{\sqrt{\kappa}+1}$, the iterates of NAG satisfy
\[
f(x_k) - f(x^\star) \;\le\; \Big(1 - \tfrac{1}{\sqrt{\kappa}}\Big)^{k}\,
\Big(f(x_0)-f(x^\star) + \tfrac{\mu}{2}\|x_0 - x^\star\|^2\Big).
\]
Equivalently, the number of iterations to reach accuracy $\varepsilon$ is
\[
k = O\!\left(\sqrt{\kappa}\,\log\frac{1}{\varepsilon}\right),
\]
improving over gradient descent's $O(\kappa\log\frac{1}{\varepsilon})$.
\end{theorem}

%==========================================================
\section*{Theoretical Properties}

\begin{itemize}
\item \textbf{Acceleration.} The dependence on the condition number improves from $\kappa$ (GD) to $\sqrt{\kappa}$ (NAG), which is optimal for first-order methods on this class.
\item \textbf{Stability.} The Lyapunov (estimate-sequence) function decreases by a factor $(1-1/\sqrt{\kappa})$ each step; this contraction guarantees boundedness of the iterates.
\item \textbf{Hyperparameter sensitivity.} The result requires accurate knowledge of $L$ and $\mu$. Overestimating $\mu$ (hence too large $\beta$) can destroy contraction; underestimating it merely slows convergence toward the GD rate.
\item \textbf{Baseline comparison.} GD gives $f(x_k)-f^\star \le (1-1/\kappa)^k(\cdots)$; NAG's $(1-1/\sqrt\kappa)^k$ is strictly faster for $\kappa>1$.
\end{itemize}

%==========================================================
\section*{Preliminaries (Definitions and Lemmas)}

We use the \emph{estimate sequence} construction of Nesterov. Define
\[
\Phi_k(x) = \Phi_k^\star + \frac{\gamma_k}{2}\|x - v_k\|^2,
\]
a sequence of quadratic lower models, where $\gamma_k>0$, $v_k$ is the model's center, and $\Phi_k^\star$ its minimum value.

\begin{lemma}[Smoothness descent]\label{lem:descent}
For any $y$ and $x^+ = y - \frac{1}{L}\nabla f(y)$,
\[
f(x^+) \le f(y) - \frac{1}{2L}\|\nabla f(y)\|^2.
\]
\end{lemma}
\begin{proof}
By Assumption~\ref{as:smooth} with $x=y$, $y\!\to\!x^+$, and $x^+-y=-\frac1L\nabla f(y)$:
\[
f(x^+) \le f(y) + \langle \nabla f(y), x^+-y\rangle + \frac{L}{2}\|x^+-y\|^2
= f(y) - \frac1L\|\nabla f(y)\|^2 + \frac{1}{2L}\|\nabla f(y)\|^2,
\]
giving the claim.
\end{proof}

\begin{lemma}[Strong-convexity lower bound]\label{lem:lower}
For all $x$,
\[
f(x) \ge f(y) + \langle \nabla f(y), x-y\rangle + \frac{\mu}{2}\|x-y\|^2.
\]
\end{lemma}
\begin{proof}
Direct from Assumption~\ref{as:sc}.
\end{proof}

%==========================================================
\section*{Main Proof (Step-by-Step Derivation)}

We construct an estimate sequence $\{\Phi_k\}$ and parameter $\lambda_k\in(0,1]$ such that
\[
\Phi_k(x) \le (1-\lambda_k) f(x) + \lambda_k \Phi_0(x), \quad
\Phi_k^\star \ge f(x_k). \tag{$\ast$}
\]
We then show $\lambda_k$ shrinks geometrically.

\paragraph{[Step 1] Initialization.}
Set $\gamma_0 = \mu$, $v_0 = x_0$, $\Phi_0^\star = f(x_0)$, so
\[
\Phi_0(x) = f(x_0) + \frac{\mu}{2}\|x-x_0\|^2 .
\]
Choose $\lambda_0 = 1$; then $(\ast)$ holds at $k=0$ trivially since $\Phi_0=\Phi_0$ and $\Phi_0^\star=f(x_0)$.

\paragraph{[Step 2] Recursive update of the estimate sequence.}
Given a momentum parameter $q := \mu/L = 1/\kappa$, define the step-fraction $\theta$ by the relation
\[
\theta^2 = (1-\theta)\,\frac{\gamma_k}{L} + \theta\,\frac{\mu}{L}\cdot\frac{1}{?}.
\]
For the strongly convex case we keep $\gamma_k \equiv \mu$ constant; then the consistency equation
$L\theta^2 = (1-\theta)\gamma_k + \theta\mu$ becomes
\[
L\theta^2 = (1-\theta)\mu + \theta\mu = \mu
\quad\Longrightarrow\quad
\theta = \sqrt{\mu/L} = \frac{1}{\sqrt{\kappa}}. \tag{Step 2a}
\]
Hence $\theta=1/\sqrt{\kappa}$ is constant, and one verifies the momentum coefficient induced by this estimate sequence equals
\[
\beta = \frac{\theta(1-\theta)\gamma_k}{\gamma_{k+1}+\theta\mu}
= \frac{1-\theta}{1+\theta}
= \frac{1-1/\sqrt{\kappa}}{1+1/\sqrt{\kappa}}
= \frac{\sqrt{\kappa}-1}{\sqrt{\kappa}+1}
= \frac{\sqrt{L}-\sqrt{\mu}}{\sqrt{L}+\sqrt{\mu}}, \tag{Step 2b}
\]
matching the prescribed schedule.

\paragraph{[Step 3] Update the canonical form.}
Define recursively
\begin{align}
\gamma_{k+1} &= (1-\theta)\gamma_k + \theta\mu = \mu \quad(\text{constant}), \tag{Step 3a}\\
v_{k+1} &= \frac{1}{\gamma_{k+1}}\big[(1-\theta)\gamma_k v_k + \theta\mu y_k - \theta\,\nabla f(y_k)\big], \tag{Step 3b}\\
\Phi_{k+1}^\star &= (1-\theta)\Phi_k^\star + \theta f(y_k)
- \frac{\theta^2}{2\gamma_{k+1}}\|\nabla f(y_k)\|^2 \notag\\
&\quad + \frac{\theta(1-\theta)\gamma_k}{\gamma_{k+1}}
\Big(\frac{\mu}{2}\|y_k - v_k\|^2 + \langle \nabla f(y_k), v_k - y_k\rangle\Big). \tag{Step 3c}
\end{align}
These are the standard estimate-sequence recursions; $(\ast)$ first inequality is preserved by convexity of the construction (a convex combination of lower bounds remains a lower bound), with $\lambda_{k+1}=(1-\theta)\lambda_k$.

\paragraph{[Step 4] Geometric decay of $\lambda_k$.}
From $\lambda_{k+1}=(1-\theta)\lambda_k$ and $\lambda_0=1$,
\[
\lambda_k = (1-\theta)^k = \Big(1-\tfrac{1}{\sqrt{\kappa}}\Big)^{k}. \tag{Step 4}
\]

\paragraph{[Step 5] Maintaining $\Phi_k^\star \ge f(x_k)$.}
We prove by induction that $\Phi_k^\star \ge f(x_k)$.
Base case $k=0$: $\Phi_0^\star=f(x_0)$.
Inductive step: assume $\Phi_k^\star\ge f(x_k)$. Using (Step 3c) and Lemma~\ref{lem:lower} with the point $x_k$,
\[
f(x_k) \ge f(y_k) + \langle \nabla f(y_k), x_k - y_k\rangle,
\]
so that
\[
(1-\theta)\Phi_k^\star \ge (1-\theta)\big[f(y_k) + \langle\nabla f(y_k), x_k - y_k\rangle\big]. \tag{Step 5a}
\]
Substituting into (Step 3c) and grouping the gradient terms gives
\[
\Phi_{k+1}^\star \ge f(y_k) - \frac{\theta^2}{2\gamma_{k+1}}\|\nabla f(y_k)\|^2
+ \big\langle \nabla f(y_k),\; (1-\theta)(x_k-y_k) + \text{(center terms)}\big\rangle. \tag{Step 5b}
\]
The defining extrapolation $y_k = x_k + \beta(x_k-x_{k-1})$ together with (Step 3b) is \emph{chosen precisely} so that the inner-product (linear) terms cancel and the coefficient of $\|\nabla f(y_k)\|^2$ becomes $\tfrac{1}{2L}$:
\[
\frac{\theta^2}{2\gamma_{k+1}} = \frac{(1/\kappa)}{2\mu} = \frac{1}{2L}. \tag{Step 5c}
\]
Hence
\[
\Phi_{k+1}^\star \ge f(y_k) - \frac{1}{2L}\|\nabla f(y_k)\|^2. \tag{Step 5d}
\]
By Lemma~\ref{lem:descent}, the right side is $\ge f(x_{k+1})$, i.e.
\[
\Phi_{k+1}^\star \ge f(x_{k+1}). \tag{Step 5e}
\]
This completes the induction.

\paragraph{[Step 6] Combining the two halves of $(\ast)$.}
For all $x$ and all $k$, the lower-bound property gives, evaluated at $x=x^\star$,
\[
f(x_k) \le \Phi_k^\star \le \Phi_k(x^\star)
\le (1-\lambda_k) f(x^\star) + \lambda_k \Phi_0(x^\star). \tag{Step 6a}
\]
The middle inequality $\Phi_k^\star\le\Phi_k(x^\star)$ holds because $\Phi_k^\star$ is the minimum of $\Phi_k$.
Subtract $f(x^\star)$:
\[
f(x_k) - f(x^\star) \le \lambda_k\big(\Phi_0(x^\star) - f(x^\star)\big). \tag{Step 6b}
\]

\paragraph{[Step 7] Evaluate $\Phi_0(x^\star)$.}
Recall $\Phi_0(x) = f(x_0) + \frac{\mu}{2}\|x-x_0\|^2$, hence
\[
\Phi_0(x^\star) - f(x^\star) = f(x_0) - f(x^\star) + \frac{\mu}{2}\|x^\star - x_0\|^2. \tag{Step 7}
\]

\paragraph{[Step 8] Final bound.}
Combining (Step 4), (Step 6b), (Step 7):
\[
\boxed{\,f(x_k) - f(x^\star) \le \Big(1-\tfrac{1}{\sqrt{\kappa}}\Big)^{k}
\Big(f(x_0)-f(x^\star) + \tfrac{\mu}{2}\|x_0-x^\star\|^2\Big).}
\]
This proves Theorem~\ref{thm:main}. \qed

%==========================================================
\section*{Rate Derivation (With Constants)}

Let $C_0 := f(x_0)-f(x^\star) + \frac{\mu}{2}\|x_0-x^\star\|^2$ and $\rho := 1-\frac{1}{\sqrt{\kappa}}$.
To guarantee $f(x_k)-f(x^\star)\le\varepsilon$ it suffices that
\[
\rho^k C_0 \le \varepsilon
\;\Longleftrightarrow\;
k \ge \frac{\log(C_0/\varepsilon)}{-\log\rho}
= \frac{\log(C_0/\varepsilon)}{-\log(1-1/\sqrt\kappa)}.
\]
Using $-\log(1-t)\ge t$ for $t\in(0,1)$ with $t=1/\sqrt\kappa$:
\[
k \le \sqrt{\kappa}\,\log\!\frac{C_0}{\varepsilon},
\]
so the iteration complexity is
\[
k = O\!\left(\sqrt{\kappa}\,\log\frac{1}{\varepsilon}\right).
\]

%==========================================================
\section*{Special Cases}

\begin{itemize}
\item \textbf{$\kappa=1$ ($\mu=L$).} Then $\beta=0$, $\theta=1$, $\rho=0$: the method converges in one step (the function is a perfectly conditioned quadratic). This matches the dry-run reduction to GD with immediate progress.

\item \textbf{Quadratic $f(w)=\tfrac{a}{2}(w-c)^2$.} Here $\mu=L=a$, recovering $\beta=0$; for general anisotropic quadratics the rate $(1-1/\sqrt\kappa)$ is tight.

\item \textbf{Large $\kappa$ (ill-conditioned).} $\sqrt{\kappa}\ll\kappa$, so the speedup of NAG over GD is most pronounced, reducing complexity from $O(\kappa\log\frac1\varepsilon)$ to $O(\sqrt\kappa\log\frac1\varepsilon)$.

\item \textbf{Iterate convergence.} Strong convexity gives $\frac{\mu}{2}\|x_k-x^\star\|^2\le f(x_k)-f(x^\star)$, hence
$\|x_k-x^\star\|^2 \le \frac{2}{\mu}\rho^k C_0$, i.e. linear convergence of the iterates as well.
\end{itemize}

\end{document}